\section*{Bab \ref{ch:g11:titration} --- Titrasi}

\begin{solution}{exo:g11:titration:1}
\omterm{def:g11:titration:titration}{Titran} adalah larutan yang konsentrasinya diketahui dengan teliti, di
dalam buret; \omterm{def:g11:titration:titration}{larutan yang dititrasi} mengandung spesies yang akan diukur,
di dalam labu. Pada \omterm{def:g11:titration:equivalence}{titik ekuivalen}, \omterm{def:g11:titration:titration}{titran} yang ditambahkan dan spesies
yang dititrasi berada dalam perbandingan yang sesuai dengan persamaan:
keduanya habis.
\end{solution}

\begin{solution}{exo:g11:titration:2}
$n = 0.0200 \times 0.0125 = \qty{2.50e-4}{mol}$.
\end{solution}

\begin{solution}{exo:g11:titration:3}
$V_E = \qty{14.3}{mL}$; garis besi(II) dimulai pada
\qty{14.3e-4}{mol}, yaitu \qty{1.43e-3}{mol}.
\end{solution}

\begin{solution}{exo:g11:titration:4}
\omterm{def:g11:titration:titration}{Titran} itu sendiri berwarna (ungu) dan produknya hampir tak berwarna:
warnanya muncul di dalam labu begitu permanganat berlebih.
\end{solution}

\begin{solution}{exo:g11:titration:5}
5: persamaannya menghabiskan 5 \omterm{def:g9:ions:ion}{ion} besi(II) per \omterm{def:g9:ions:ion}{ion} permanganat, dan
pada \omterm{def:g11:titration:equivalence}{titik ekuivalen} \omterm{def:g7:chemical-reactions:reactant}{pereaksi-pereaksinya} telah dipertemukan dalam
perbandingan itu.
\end{solution}

\begin{solution}{exo:g11:titration:6}
$n(\ce{MnO4^-}) = 0.0200 \times 0.00840 = \qty{1.68e-4}{mol}$;
$n(\ce{Fe^{2+}}) = 5 \times \num{1.68e-4} = \qty{8.40e-4}{mol}$;
$c = \num{8.40e-4} / 0.0100 = \qty{0.0840}{mol/L}$;
$C_m = 0.0840 \times 55.8 = \qty{4.69}{g/L}$.
\end{solution}

\begin{solution}{exo:g11:titration:7}
Perbandingan 1:1: $n = 0.0100 \times 0.0114 = \qty{1.14e-4}{mol}$ dalam
\qty{10.0}{mL}, jadi \qty{1.14e-3}{mol} dalam \qty{100.0}{mL}, yaitu
$\num{1.14e-3} \times 176.0 = \qty{0.201}{g}$: sekitar \qty{200}{mg}.
\end{solution}

\begin{solution}{exo:g11:titration:8}
Setiap tetes tidak akan langsung bereaksi: warna \omterm{def:g11:titration:titration}{titran} akan bertahan
sebelum \omterm{def:g11:titration:equivalence}{titik ekuivalen}, dan akhirnya akan terlihat terlalu cepat (atau
orang harus menunggu beberapa menit setelah setiap tetes).
\end{solution}

\begin{solution}{exo:g11:titration:9}
Terlalu besar: ia telah melewati \omterm{def:g11:titration:equivalence}{titik ekuivalen}, menambahkan
permanganat berlebih. Mendekati \omterm{def:g11:titration:equivalence}{titik ekuivalen} seharusnya ia
menambahkan \omterm{def:g11:titration:titration}{titran} tetes demi tetes, sambil menggoyang labu, lalu
berhenti pada warna merah muda samar pertama yang bertahan.
\end{solution}

\begin{solution}{exo:g11:titration:10}
$n(\ce{Fe^{2+}}) = 5 \times 0.0200 \times 0.0130 = \qty{1.30e-3}{mol}$
dalam \qty{10.0}{mL}: $c = \qty{0.130}{mol/L}$ untuk larutan encer, dan
$10 \times 0.130 = \qty{1.30}{mol/L}$ untuk larutan pekatnya.
\end{solution}

\begin{solution}{exo:g11:titration:11}
$0.0630 \times 55.8 = \qty{3.52}{g}$.
\end{solution}

\begin{solution}{exo:g11:titration:12}
$0.1 / 7.2 \approx \qty{1.4}{\percent}$, dibandingkan
\qty{0.7}{\percent} untuk \qty{14.3}{mL}. Titrasilah volume larutan yang
lebih besar (atau sampel yang lebih pekat), atau pakailah \omterm{def:g11:titration:titration}{titran} yang
lebih encer.
\end{solution}

\begin{solution}{exo:g11:titration:13}
$n(\ce{MnO4^-}) = 0.0200 \times 0.0176 = \qty{3.52e-4}{mol}$;
$n(\ce{H2O2}) = \frac{5}{2} \times \num{3.52e-4} = \qty{8.80e-4}{mol}$
dalam \qty{10.0}{mL}: \qty{0.0880}{mol/L} dalam larutan encer,
\qty{0.880}{mol/L} dalam disinfektan, yaitu
$0.880 \times 34.0 = \qty{29.9}{g/L}$ (larutan sekitar 3 \%).
\end{solution}

\begin{solution}{exo:g11:titration:14}
$n(\ce{MnO4^-}) = \num{1.4e-3} / 5 = \qty{2.8e-4}{mol}$, yang harus
dialirkan dalam sekitar \qty{0.015}{L}: $c \approx \qty{0.019}{mol/L}$,
jadi larutan \qty{0.0200}{mol/L}. Jika sepuluh kali lebih pekat, $V_E$
akan sekitar \qty{1.4}{mL}: kesalahan buret akan mencapai sekitar
\qty{7}{\percent}.
\end{solution}

\begin{solution}{exo:g11:titration:15}
$8 \times 0.0200 \times 0.0143 = \qty{2.29e-3}{mol}$ \omterm{def:g9:ions:ion}{ion} hidrogen.
\qty{10}{mL} pada \qty{1.0}{mol/L} mengandung \qty{1.0e-2}{mol}: cukup,
lebih dari empat kali lipat. Tanpa asam, reaksi dalam bab ini tidak dapat
berlangsung seperti yang tertulis (\omterm{def:g9:ions:ion}{ion} hidrogen adalah \omterm{def:g7:chemical-reactions:reactant}{pereaksinya}), dan
permanganat akan bereaksi secara lain.
\end{solution}

\begin{solution}{pb:g11:titration:1}
\textbf{1.} \ce{MnO4^-}/\ce{Mn^{2+}} dan \ce{Fe^{3+}}/\ce{Fe^{2+}}.
\omterm{def:g11:redox:oxidant}{Oksidatornya} \omterm{def:g9:ions:ion}{ion} permanganat, \omterm{def:g11:redox:oxidant}{reduktornya} \omterm{def:g9:ions:ion}{ion} besi(II).

\textbf{2.} \ce{MnO4^- + 8H+ + 5e- -> Mn^{2+} + 4H2O};
\ce{Fe^{2+} -> Fe^{3+} + e-}.

\textbf{3.} Yang kedua dikalikan 5:
\ce{MnO4^- + 8H+ + 5Fe^{2+} -> Mn^{2+} + 5Fe^{3+} + 4H2O}. Muatan di
kiri $-1 + 8 + 10 = +17$; di kanan $+2 + 15 = +17$.

\textbf{4.} \omterm{def:g9:ions:ion}{Ion} hidrogen adalah \omterm{def:g7:chemical-reactions:reactant}{pereaksi}: reaksinya memerlukan \omterm{def:g9:acids-bases-ph:acidic}{larutan
asam}.

\textbf{5.} Reaksinya tuntas, cepat, dan satu-satunya reaksi permanganat
di dalam labu; permanganat yang ungu habis terpakai sebelum \omterm{def:g11:titration:equivalence}{titik
ekuivalen} dan mewarnai labu menjadi merah muda sesudahnya.

\textbf{6.} Di dalam buret; dibaca pada dasar meniskus, mata sejajar,
sebelum dan sesudah, sampai \qty{0.05}{mL}.

\textbf{7.} Kuning pucat (\omterm{def:g9:ions:ion}{ion-ion} besi); setiap tetes ungu langsung
dihabiskan oleh \omterm{def:g9:ions:ion}{ion-ion} besi(II) yang masih ada.

\textbf{8.} $0.0200 \times 0.0143 = \qty{2.86e-4}{mol}$.

\textbf{9.} $n(\ce{Fe^{2+}}) = 5 \times \num{2.86e-4} =
\qty{1.43e-3}{mol}$.

\textbf{10.} $\num{1.43e-3} \times 55.8 = \qty{0.0798}{g} =
\qty{79.8}{mg}$.

\textbf{11.} $(80 - 79.8)/80 \approx \qty{0.3}{\percent}$.

\textbf{12.} $\qty{0.1}{mL}$ dari \qty{14.3}{mL} adalah
\qty{0.7}{\percent}, sekitar \qty{0.6}{mg}: hasilnya, $79.8 \pm 0.6$ mg,
sesuai dengan labelnya.

\textbf{13.} $14.1$: $5 \times 0.0200 \times 0.0141 \times 55.8 =
\qty{78.7}{mg}$; $14.4$: \qty{80.4}{mg}.

\textbf{14.} $(79.8 + 78.7 + 80.4)/3 = \qty{79.6}{mg}$.

\textbf{15.} Tablet-tablet itu sendiri sedikit berbeda dalam kandungan
besinya, sekitar \qty{1}{\percent}.

\textbf{16.} Tidak: $V_E$ akan sekitar \qty{2.9}{mL}, dan kesalahan
buretnya \qty{3.5}{\percent}, lima kali lebih buruk.

\textbf{17.} Vitamin C juga akan dioksidasi oleh permanganat: reaksinya
tidak lagi tunggal, $V_E$ akan terlalu besar, dan besinya akan
dihitung terlalu tinggi.

\textbf{18.} $\num{1.43e-3} \times 151.9 = \qty{0.217}{g}$, sekitar
\qty{217}{mg}.

\textbf{19.} $\num{1.43e-3} \times \num{6.02e23} = \num{8.6e20}$ \omterm{def:g9:ions:ion}{ion}.

\textbf{20.} Sekitar \qty{80}{mg} besi per tablet: labelnya jujur.
\end{solution}
